\src{109}
các tập đơn hình, tức các hàm tử phản biến $K$ từ $\mathcal S$ vào
$(\mathrm{Ens})$, biến các phép đẩy ra thành các tích thớ.

Do đó, có thể xem các phạm trù như những tập bán đơn hình đặc biệt, và các
groupoid như những tập đơn hình đặc biệt, dĩ nhiên với điều kiện phải lập luận
«sai khác một đẳng cấu», như bắt buộc khi diễn giải một số cấu trúc bằng những
cấu trúc khác. Phương pháp thông thường quy về trường hợp lý thuyết tập hợp
khi đó kéo theo:

\result{HỆ QUẢ}{4.2} Mệnh đề trên vẫn đúng khi thay các phạm trù, các
groupoid và các tập đơn hình bằng các $\mathcal C$-phạm trù,
$\mathcal C$-groupoid và các $\mathcal C$-vật đơn hình, với điều kiện các tích
thớ tồn tại trong $\mathcal C$.

Vật bán đơn hình $K$ trong $\mathcal C$ liên kết với một phạm trù
$(X,R,\ldots)$ trong $\mathcal C$ có thể được viết tường minh bằng cách diễn
giải thành phần $K_n=K(\Delta_n)$ của $K$ là tích thớ lặp $(n+1)$ lần của
$(R,p_1)$ trên $X$, hoặc tốt hơn, bằng công thức truy hồi
\[
  K_0=R,\qquad
  K_n=(K_{n-1},p_n^{(n-1)})\times_X(R,p_1),
\]
trong đó các $p_i^{(n-1)}$ ($0\leq i\leq n-1$) là các phép chiếu tự nhiên từ
$K_{n-1}$ xuống $X$ (cũng được định nghĩa bằng truy hồi). Theo cách này,
$p_1,p_2,\pi,D,s$ được diễn giải là các phép toán đơn hình lần lượt tương ứng
với các cấu xạ sau trong $\mathcal S$: mặt $0$ của $\Delta_1$, mặt $1$ của
$\Delta_1$, mặt $(0,2)$ của $\Delta_2$, phép suy biến
$\Delta_1\to\Delta_0$, và phép đối xứng của $\Delta_1$. Mọi phép toán bán đơn
hình khác (tương ứng, đơn hình) đều được suy ra một cách hình thức từ bốn
phép toán trước (tương ứng, năm phép toán) bằng phép hợp thành và các tích thớ.

Nay gọi \emph{tiền quan hệ tương đương} trên một vật $X$ của một phạm trù là
việc cho một groupoid mà vật của các vật (nếu có thể nói vậy) là $X$. Vì thế,
dữ liệu ấy bao gồm, ngoài những thứ khác, một vật $R$ và hai cấu xạ
\[
  p_1,p_2:R\rightrightarrows X.
\]
Tuy nhiên, lưu ý rằng riêng các dữ liệu này không xác định cấu trúc đang xét,
khác với trường hợp các cặp tương đương. Trong exposé này, ta quan tâm đến
khái niệm ấy nhằm thu được các tiêu chuẩn để có thể lấy thương, tức là lập đối
hạt nhân của cặp $(p_1,p_2)$. Vì vậy, phát biểu của bài toán không dùng đến
các dữ liệu bổ sung vốn có trong một groupoid. Trong việc chứng minh các kết
quả sẽ nêu, ta sẽ
